% 第 3 册（大学化学，第二学年）照片与 AI 插图的图片致谢。
% 机器可读的许可证记录保存在
% images/book3/CREDITS.md 中；两份文件须保持同步。第 3 册的撰写者
% 负责本文件：每收到一张照片，就在 itemize 列表中（放在
% 第一段之后）添加一个 \item。
\chapter*{图片致谢}

本书中的插图均为原创。照片（如有使用）依各自的自由许可证使用，
谨向作者致谢；
每张照片的完整来源链接与许可证网址均列于本书代码库中的
\texttt{images/book3/CREDITS.md} 文件。
危险象形图采用的是全球化学品统一分类和标签制度（GHS）的图形，
\texttt{ghsystem} 宏包（\TeX\ Live）所绘制的版本。

\bigskip
除照片和原创插图外，部分图形为
AI 生成的插图，由 OpenAI 图像生成功能产出，提示词
由本书编辑撰写，并在收入前经过化学
准确性审核。每张生成图像的完整提示词均记录于
\texttt{images/book3/ai/PROMPTS.md} 中（位于代码库内）。

\bigskip
\noindent 照片，按章节列出：
\begin{itemize}\small
  \item 第~1 章：Germain Henri Hess，I.~I. Reimers 绘制的肖像（1837--1838），
        CC BY-SA 4.0；Marcellin Berthelot，作者不详，公有领域
        （均来自 维基共享资源）。
  \item 第~2 章：Josiah Willard Gibbs，摄影者不详，公有领域
        （维基共享资源）。
  \item 第~3 章：François-Marie Raoult，摄影者不详，公有领域
        （维基共享资源）。
  \item 第~4 章：Jacobus Henricus van 't Hoff，Nicola Perscheid 作（1904），
        公有领域；Henry Le Chatelier，作者不详，CC BY-SA 4.0；
        Fritz Haber 与 Carl Bosch，诺贝尔基金会，公有领域（均
        来自 维基共享资源）。
  \item 第~6 章：Duisburg-Nord 的高炉，Dietmar Rabich 作，CC BY-SA
        4.0；钢轨的铝热焊接，PetrS. 作，CC BY-SA 3.0（均
        来自 维基共享资源）。
  \item 第~7 章：炼油厂的精馏塔，CEphoto、Uwe Aranas 作，
        CC BY-SA 3.0（维基共享资源）。
  \item 第~9 章：Michael Faraday，Thomas Phillips 作（1842），公有领域
        （维基共享资源）。
  \item 第~10 章：Jaroslav Heyrovsk\'y，J.~Heyrovsk\'y 物理化学研究所档案，
        CC BY-SA 3.0（维基共享资源）。
  \item 第~11 章：伏打电堆，GuidoB 作，CC BY-SA 3.0；Charles Martin Hall
        与 Paul H\'eroult，作者不详，公有领域（均 Wikimedia
        Commons）。
  \item 第~12 章：铜像，Elcobbola 作，公有领域；牺牲阳极，
        Zwergelstern 作，CC BY-SA 3.0（均来自 维基共享资源）。
  \item 第~13 章：Erwin Schr\"odinger，诺贝尔基金会（1933），公有领域
        （维基共享资源）。
  \item 第~14 章：磁铁中的液氧，Pieter Kuiper 作，公有领域；
        Robert Mulliken 与 Friedrich Hund，GFHund 作，CC BY 3.0（均 Wikimedia
        Commons）。
  \item 第~16 章：August Kekul\'e（1873），作者不详，公有领域
        （维基共享资源）。
  \item 第~18 章：过渡金属溶液，Benjah-bmm27 作，公有领域；Alfred
        Werner，Les Prix Nobel（1914），公有领域（均来自 维基共享资源）。
  \item 第~19 章：红宝石晶体，公有领域；祖母绿，Didier Descouens 作，
        CC BY-SA 3.0；硫酸铜晶体，Stephanb 作，CC BY-SA 3.0（均
        来自 维基共享资源）。
  \item 第~20 章：Akira Suzuki、Ei-ichi Negishi 与 Richard Heck，BloodIce 作，CC BY-SA
        4.0（维基共享资源）。
  \item 第~22 章：Charles Friedel（1890 年代）与 James Mason Crafts（\textit{Popular
        Science Monthly}，1900），公有领域（维基共享资源）。
  \item 第~24 章：Michel Eug\`ene Chevreul，Nicolas Eustache Maurin 作，公有领域
        （维基共享资源）。
  \item 第~25 章：Alexander Borodin，作者不详，公有领域；Ludwig Claisen，
        Veronica.villagrasa 作，CC BY-SA 4.0（均来自 维基共享资源）。
  \item 第~26 章：Robert Robinson，作者不详，公有领域（维基共享资源）。
  \item 第~28 章：Elias James Corey，Trvthchem 作，CC0（维基共享资源）。
  \item 第~29 章：实验室中的 Wallace Carothers，摄影者不详，公有领域
        （维基共享资源）。
  \item 第~30 章：Emil Fischer，Atelier Victoria，柏林，约 1895 年，公有领域（Wikimedia
        Commons）。
  \item 第~31 章：Mikhail Tsvet，作者不详，公有领域（维基共享资源）。
  \item 第~32 章：Francis William Aston（1922），作者不详，公有领域；Gustav Kirchhoff，
        Robert Bunsen 与 Henry Roscoe（1862），公有领域（均来自 维基共享资源）。
  \item 第~34 章：William Sealy Gosset（1908），公有领域（维基共享资源）。
%% 下一章在此
\end{itemize}
\clearpage
